\section{Original coordinates and the retained cover}
The original source coordinate triple is $(x,y,w)$ and its target is
$(a,b,c)$. Its polynomial is
\begin{align}
F_1&=(1+xy)^3w+y^2(1+xy)(4+3xy),\notag\\
F_2&=y+3x(1+xy)^2w+3xy^2(4+3xy),\notag\\
F_3&=2x-3x^2y-x^3w. \label{hc:F}
\end{align}
All coordinate spaces here are complex unless a real locus is stated.
The original $x\ne0$ inverse chart is
\begin{gather}
P_{a,b,c}(r)=cr^3-2r^2+br-2a,\qquad
\alpha=\tfrac12P_r(r),\notag\\
(x,y,w)=\bigl(\alpha^{-1},r-\alpha,
                  5\alpha^2-3r\alpha-c\alpha^3\bigr),\qquad \alpha\ne0.
\label{hc:inverse}
\end{gather}
Here is the complete verification needed for this chart. At a point
with $x\ne0$, set $\alpha=1/x$ and $r=y+1/x$; solving $F_3=c$
gives the third coordinate in \eqref{hc:inverse}. Substitution into
the other original components gives
\[
F_2=4r+2\alpha-3cr^2,\qquad F_1=r^2+r\alpha-cr^3.
\]
The first equation equals $b$ precisely when $2\alpha=P_r(r)$;
after that substitution $2(F_1-a)=P(r)$. This proves both directions
and uniqueness of the root label. A multiple root gives no finite
point in this chart. Directly, $F(0,y,w)=(w+4y^2,y,0)$, so $c=0$
has the additional inverse point $(0,b,a-4b^2)$ and $c\ne0$ has none.

On $b=c=0$ the root cover is $a=-r^2$, with $\alpha=-2r$.
The arithmetic coordinate is the specific further map
\begin{equation}
a=2s,\qquad s=-r^2/2. \label{hc:cover}
\end{equation}
Its unramified domain is $\C_r^*=\C_r\setminus\{0\}$ over
$\C_s^*$. The point $r=0$ is the ramification point of the completed
cover and is excluded by $\alpha\ne0$ from the original inverse chart.
These are explicit restrictions of the original map.

The constant-coefficient heat evolution in the original root variable
is a distinct, explicitly retained map of this very polynomial:
\[
P(\tau,r)=P_{a_0+2\tau,b_0+6c\tau,c}(r),\qquad
P_\tau=6cr-4=P_{rr}.
\]
Here $\tau$ is forward heat time. The explicit reversal
$\widetilde P(t,r)=P(-t,r)$ satisfies
$\widetilde P_t=-\widetilde P_{rr}$. This is the backward sign convention;
the arithmetic Newman generator below acts in its specified $s$ coordinate.
The arithmetic pullback below transports differentiation in $s=a/2$,
which must also carry the derivative of $r$ with respect to $s$.
Both operators will therefore be displayed with their exact coordinates.

For an open set $V\subset\C_s^*$, set
$\widetilde V=\{r:-r^2/2\in V\}$. Every holomorphic function on
$\widetilde V$ has a unique decomposition
\begin{equation}
f(r)=q_0(s)+r q_1(s),\qquad q_0,q_1\in\mathcal O(V).
\label{hc:even-odd}
\end{equation}
Indeed its even part is $(f(r)+f(-r))/2$ and its odd coefficient is
$(f(r)-f(-r))/(2r)$. Both are invariant under $r\mapsto-r$ and
therefore define holomorphic functions of $s$ on local inverse
charts, agreeing on their overlaps. These formulas also prove uniqueness.
They define an isomorphism $T_V:\mathcal O(V)^2\to\mathcal O(\widetilde V)$.
The same decomposition holds over $V$ containing zero: power series
show that the odd numerator divided by $r$ extends holomorphically
and that each even power series is a power series in $s=-r^2/2$.
The derivative operator below, however, is first defined only over
$V\subset\C^*$.

\section{Connection, square, and branch coupling}
Write a section as the column $q=(q_0,q_1)^{\mathsf T}$ in the
ordered basis $(1,r)$. On the unramified cover,
$ds/dr=-r$, so differentiation in the original spectral coordinate is
$\partial_s^{\rm lift}=-r^{-1}\partial_r$. In particular
$\partial_s r=-1/r=r/(2s)$. Consequently
\begin{equation}
D=\dd+A(s),\quad A(s)=\begin{pmatrix}0&0\\0&1/(2s)\end{pmatrix},
\qquad T_VD=\partial_s^{\rm lift}T_V.
\label{hc:D}
\end{equation}
This is a connection: for $g\in\mathcal O(V)$,
$D(gq)=g'q+gDq$. Its curvature in this one complex coordinate is zero;
the nontrivial information is its ramification and monodromy.

\begin{proposition}[The full second derivative and its commutators]
On $\mathcal O(V)^2$ one has
\begin{align}
D^2&=\begin{pmatrix}
\partial_s^2&0\\0&\partial_s^2+s^{-1}\partial_s-1/(4s^2)
\end{pmatrix}, \label{hc:D2}\\
R&=\begin{pmatrix}0&-2s\\1&0\end{pmatrix},\quad
R^2=-2sI,\quad [D,R]=-R^{-1}, \label{hc:R}\\
[D^2,R]&=-2R^{-1}D-R^{-3}. \label{hc:D2R}
\end{align}
Here $R$ is multiplication by the original root $r$; commutators
include differentiation of the matrix coefficients.
\end{proposition}
\begin{proof}
Expanding $(\partial_s+A)^2$ gives
$\partial_s^2+2A\partial_s+A'+A^2$.
Its lower diagonal constant is
$-1/(2s^2)+1/(4s^2)=-1/(4s^2)$, proving \eqref{hc:D2}.
Multiplication $r(q_0+rq_1)=-2sq_1+rq_0$ proves $R$ and $R^2$.
Matrix multiplication then gives
\[
R'+AR-RA=\begin{pmatrix}0&-1\\1/(2s)&0\end{pmatrix}=-R^{-1}.
\]
Differentiating $RR^{-1}=I$ covariantly gives
$[D,R^{-1}]=-R^{-1}[D,R]R^{-1}=R^{-3}$.
Finally $[D^2,R]=D[D,R]+[D,R]D=-2R^{-1}D-R^{-3}$.
This proves every term, including the zeroth-order branch term.
\end{proof}

The sheet involution is $J=\operatorname{diag}(1,-1)$; it commutes
with $D,D^2$ and anticommutes with $R$.
The trace is $(q_0,q_1)\mapsto2q_0$, with kernel the odd summand,
and its half has the section $q_0\mapsto(q_0,0)$.
$D$ preserves the two summands on the punctured base, whereas $R$
exchanges them and couples the odd coefficient into $-2sq_1$.
This describes exactly the information lost by tracing.

On a simply connected $V\subset\C^*$ choose a holomorphic root
$r(s)$ satisfying $r(s)^2=-2s$. The invertible matrix
$G=\operatorname{diag}(1,r(s))$ obeys
$D=G^{-1}\partial_sG$.
Horizontal coefficients satisfy $q_0'=0$ and
$q_1'+q_1/(2s)=0$, hence $(q_0,q_1)=(c_0,c_1/r(s))$.
Continuation once around zero changes the sign of the latter
coefficient. Thus the monodromy is exactly
$\operatorname{diag}(1,-1)$; no single-valued global choice of
$r(s)$ on $\C^*$ is implied by this local gauge formula.

\section{The actual arithmetic heat function}
We use the convention of \cite[Section 1, equations (1)--(4)]{rodgersTao2020}:
\begin{align}
\Phi(u)&=\sum_{n\ge1}(2\pi^2n^4e^{9u}-3\pi n^2e^{5u})
                       e^{-\pi n^2e^{4u}},\notag\\
H_t(z)&=\int_0^\infty e^{tu^2}\Phi(u)\cos(zu)\,du,\notag\\
Z_t(s)&=8H_t(2s),\qquad
Z_0(s)=\Xi(s)=\xi(1/2+is),\notag\\
\xi(\sigma)&=\tfrac12\sigma(\sigma-1)\pi^{-\sigma/2}
                     \Gamma(\sigma/2)\zeta(\sigma).
\label{hc:Z}
\end{align}
The Mellin and spectral coordinates are related by
$\sigma=1/2+is$, with inverse $s=-i(\sigma-1/2)$.
They are not silently identified.

We rederive the integral identity and its analytic bounds. Set
\[
h(x)=\frac\pi2x^2(2\pi x^2-3)e^{-\pi x^2},\quad
S(\lambda)=\sum_{n\ge1}h(n\lambda),\quad
k(v)=e^{v/2}S(e^v).
\]
For the Fourier convention $\widehat h(q)=\int_\R h(x)e^{-2\pi ixq}dx$,
the Gaussian transform and its second and fourth derivatives give
\begin{align*}
\widehat{x^2e^{-\pi x^2}}&=(1/(2\pi)-q^2)e^{-\pi q^2},\\
\widehat{x^4e^{-\pi x^2}}&=(q^4-3q^2/\pi+3/(4\pi^2))e^{-\pi q^2}.
\end{align*}
Multiplying by $-3\pi/2$ and $\pi^2$ respectively gives $\widehat h=h$.
Since $h(0)=0$, Poisson summation gives
$S(\lambda)=\lambda^{-1}S(1/\lambda)$ and $k(v)=k(-v)$.
For $\lambda\ge1$,
$|S(\lambda)|\le C\lambda^4e^{-\pi\lambda^2}$, because the
remaining series is bounded by a constant times
$\sum_{n\ge1}n^4e^{-\pi(n^2-1)}<\infty$.
The symmetry gives the corresponding bound as $\lambda\downarrow0$.
Each fixed number of $v$ derivatives introduces only finitely many
additional powers of $ne^v$. It follows, with constants depending
on the derivative order, that
\begin{equation}
|k^{(j)}(v)|\le C_j e^{C_j|v|}e^{-\pi e^{2|v|}}.
\label{hc:kernel-bound}
\end{equation}
Termwise differentiation on compact sets is justified by the same
summable bounds. On the negative half-line differentiate the proved
evenness identity. This proves \eqref{hc:kernel-bound} on both ends.

For $\Re\sigma>1$, absolute integrability permits summation inside
the Mellin integral. The substitution $q=\pi x^2$ gives
\begin{align*}
\int_0^\infty h(x)x^{\sigma-1}dx
 &=\pi^{-\sigma/2}
   \left[\tfrac12\Gamma(\sigma/2+2)-\tfrac34\Gamma(\sigma/2+1)\right]\\
 &=\frac{\sigma(\sigma-1)}8\pi^{-\sigma/2}\Gamma(\sigma/2),\\
\int_0^\infty S(\lambda)\lambda^{\sigma-1}d\lambda&=\xi(\sigma)/4.
\end{align*}
The two-ended bounds make the last integral entire in $\sigma$,
so the equality extends by analytic continuation. Substituting
$\lambda=e^v$ and $\sigma=1/2+is$ proves
\begin{equation}
Z_t(s)=4\int_\R e^{tv^2/4}k(v)e^{isv}\,dv.
\label{hc:Fourier}
\end{equation}
Indeed direct expansion also gives $\Phi(u)=2k(2u)$; the substitution
$v=2u$ and the evenness of $k$ carry the stated $H_t$ integral
to \eqref{hc:Fourier}, retaining the factors $2,4,8$.
The bound \eqref{hc:kernel-bound} dominates every fixed derivative
uniformly on compact complex $(t,s)$ sets. Thus $Z$ is jointly entire
and differentiation under its integral proves exactly
\begin{equation}
\partial_t Z_t=-\tfrac14\partial_s^2 Z_t.
\label{hc:heat}
\end{equation}
This supplies the actual initial datum and the actual time orientation.

\section{An explicit domain for the heat evolution}
The differential expression alone does not specify a heat evolution
on every entire-function space. We now give a precise invariant domain
containing the actual kernel, without assigning it to the topology of
another spectral construction.

\begin{definition}
Let $\E$ consist of $\phi\in C^\infty(\R,\C)$ for which
\[
p_{A,B,m,j}(\phi)=\int_\R |\phi^{(j)}(v)|(1+|v|)^m
                          e^{Av^2+B|v|}\,dv<\infty
\]
for every nonnegative integer $A,B,m,j$. Give it these seminorms.
Define
\[
T\phi(s)=4\int_\R\phi(v)e^{isv}\,dv,\qquad \A=T\E,
\]
and equip $\A$ with the transported topology.
\end{definition}

The map $T$ is injective. To give the needed uniqueness argument,
let $T\phi$ vanish on the real line and let
$g_\epsilon(v)=(4\pi\epsilon)^{-1/2}e^{-v^2/(4\epsilon)}$.
Its Gaussian integral representation is
$g_\epsilon(v)=(2\pi)^{-1}\int_\R e^{-\epsilon s^2}e^{isv}ds$.
Fubini is valid because $\phi\in L^1$ and the Gaussian is integrable;
it expresses $\phi*g_\epsilon$ as a Gaussian-weighted integral of the
vanishing Fourier transform. Thus $\phi*g_\epsilon=0$.
As $\epsilon\downarrow0$, these convolutions converge to $\phi$ in
$L^1$: translations are continuous in $L^1$, first for continuous
compactly supported functions by uniform continuity, then for all
$L^1$ functions by approximation; splitting the Gaussian integral
into a small translation ball and its vanishing tail proves the claim.
Consequently $\phi=0$ almost everywhere and, being continuous, everywhere.
This also justifies the transported topology without an unspecified
choice of representative.

\begin{proposition}[Heat evolution with its inverse]
For every complex $t$, multiplication
$E_t\phi(v)=e^{tv^2/4}\phi(v)$ is a continuous automorphism of $\E$.
The operators $U_t=TE_tT^{-1}$ on $\A$ satisfy
\begin{gather}
U_tU_u=U_{t+u},\qquad U_t^{-1}=U_{-t},\qquad
\partial_tU_t=-\tfrac14\partial_s^2U_t,\notag\\
U_t\Xi=Z_t,\qquad U_t\partial_s=\partial_s U_t,\qquad
U_tM_sU_{-t}=M_s-\tfrac t2\partial_s. \label{hc:conjugation}
\end{gather}
Here $M_s$ is multiplication by the unchanged spectral coordinate.
\end{proposition}
\begin{proof}
Every $v$ derivative of $e^{tv^2/4}\phi$ is a finite sum of
$e^{tv^2/4}$ times a polynomial in $v$ times a derivative of $\phi$.
Each target seminorm is bounded by finitely many source seminorms
with larger $A,m$; on compact $t$ sets these bounds are uniform.
The inverse multiplier is $e^{-tv^2/4}$. Differentiation in $t$
inserts $v^2/4$, while two $s$ derivatives in the transform insert
$-v^2$, proving the generator and the group identities.
The kernel bound proves $k\in\E$, so $U_t\Xi=Z_t$ follows from
\eqref{hc:Fourier}.

Integration by parts gives $M_sT\phi=T(i\phi')$.
For completeness the boundary terms vanish: each derivative of
$e^{isv}\phi$ is integrable on a horizontal real-$v$ line and the
function itself is integrable, so its limits at both ends exist and
are zero. The same argument applies after exponential weights.
Thus multiplication by $s$ preserves $\A$; differentiation does also,
since $\partial_sT\phi=T(iv\phi)$.
Finally
\[
e^{tv^2/4}i\partial_v(e^{-tv^2/4}\phi)
   =i\phi'-\tfrac t2iv\phi.
\]
Applying $T$ proves the last identity in \eqref{hc:conjugation}.
All displayed operators have the domains just specified.
\end{proof}

\section{Transport of the moving arithmetic relation}
Put $B_t=M_s-(t/2)\partial_s$. For every polynomial $p$,
the automorphism calculation proves
\begin{equation}
U_t\bigl(p(s)\Xi(s)\bigr)=p(B_t)Z_t(s).
\label{hc:relation}
\end{equation}
The right side is the polynomial in the actual differential operator,
with composition order retained. For example the first two nonconstant
relations are
\begin{align*}
B_tZ_t&=sZ_t-\tfrac t2Z_t',\\
B_t^2Z_t&=(s^2-t/2)Z_t-tsZ_t'+\tfrac{t^2}4Z_t''.
\end{align*}
These follow from $\partial_s(sq)=q+s\partial_s q$.
In particular, an evolving polynomial multiple is not obtained by
silently leaving multiplication by $s$ fixed.

One can take a closure in the topology that has actually been defined:
let $N_0=\overline{\C[s]\Xi}^{\A}$ and $N_t=U_tN_0$.
Continuity and invertibility give
$N_t=\overline{\{p(B_t)Z_t:p\in\C[s]\}}^{\A}$.
$M_s$ preserves $N_0$, since it is continuous and preserves its dense
generating subspace; hence $B_t$ preserves $N_t$.
The map $[f]\mapsto[U_tf]$ is a well-defined continuous isomorphism
$\A/N_0\to\A/N_t$ with inverse induced by $U_{-t}$ and intertwines
the induced $M_s$ with the induced $B_t$.
This proof applies to these explicit spaces. It makes no identification
with the topology or completion of the Connes quotient
\cite[Section 3.6]{ccmProlate2024}.

There is a second exact description, local at the moving divisor.
Let $Z=Z_t(s)$ and $\mathcal P=\partial_t+\tfrac14\partial_s^2$.
For every holomorphic $g(t,s)$ the complete product calculation is
\begin{equation}
\mathcal P(Zg)=Z\mathcal Pg+\tfrac12Z_sg_s. \label{hc:ideal-defect}
\end{equation}
The term $\mathcal PZ$ is zero by \eqref{hc:heat}; the mixed product
term is retained. On the complement of $Z=0$, set $b=Z_s/Z$.
Direct differentiation of $f/Z$ gives the exact conjugation
\begin{equation}
Z\mathcal P Z^{-1}
=\partial_t+\tfrac14\partial_s^2-\tfrac12 b\partial_s+\tfrac12 b^2.
\label{hc:gauge}
\end{equation}
To check the potential term before canceling any expression, it is
$\tfrac12(Z_s/Z)^2-\tfrac14Z_{ss}/Z-Z_t/Z$;
the last two terms cancel by the established heat equation.
Thus the operator in \eqref{hc:gauge} sends $Zg$ exactly to
$Z\mathcal Pg$. Its coefficients are meromorphic at zeros of $Z$,
but its restriction to the ideal of holomorphic multiples of $Z$
has the displayed holomorphic value. Division by $Z$ on this ideal
is well-defined and unique by analytic continuation. There is no
claim that division is holomorphic on arbitrary sections at a zero.

The identical calculation on the two-component module replaces
$\partial_s$ by $D$:
\begin{equation}
Z(\partial_t+\tfrac14D^2)Z^{-1}
=\partial_t+\tfrac14D^2-\tfrac12 bD+\tfrac12b^2I.
\label{hc:module-gauge}
\end{equation}
Indeed $D(Zq)=Z_s q+ZDq$, so the product expansion is exactly the
scalar one, retaining $A(s)$ in $D$. This gives the full moving
arithmetic relation on the retained cover without assuming a derivative
acts on a frozen quotient.

\section{The exact branch-sensitive heat defect}
The actual pulled arithmetic solution is $q(t,s)=(Z_t(s),0)$.
Equation \eqref{hc:D2} proves
$(\partial_t+\tfrac14D^2)q=0$.
Its function on the cover is $f(t,r)=Z_t(-r^2/2)$.
By differentiating $s=-r^2/2$ twice, the scalar generator is
\begin{equation}
L_r=-\frac1{4r^2}\partial_r^2+\frac1{4r^3}\partial_r,
\qquad \partial_t f=L_rf\quad(r\ne0).
\label{hc:Lr}
\end{equation}
This retains the original spectral diffusion; it does not substitute
the constant-coefficient $r$-heat equation of the fibre polynomial.

The branch observable $rZ_t$ is the column $Rq=(0,Z_t)$.
From the already proved commutator its complete defect is
\begin{align}
(\partial_t+\tfrac14D^2)(Rq)
&=-\tfrac12R^{-1}Dq-\tfrac14R^{-3}q\notag\\
&=\begin{pmatrix}0\\Z_t'/(4s)-Z_t/(16s^2)\end{pmatrix}.
\label{hc:odd-defect}
\end{align}
Equivalently, as a function on the cover it is
$-Z_t'/(2r)-Z_t/(4r^3)$.
These terms are the exact dynamics of the retained odd observable,
not a reason to discard that observable or declare it unrelated.
For real $t$, $Z_t(0)>0$ (proved from the integral in
\eqref{eq:zt-imaginary-axis-positive}), so its leading branch pole is nonzero. The original
arithmetic solution itself remains entire at the branch point.

\section{Keeping the other original fibre coefficients}
The preceding branch at $s=0$ belongs to the particular slice $b=c=0$.
We now retain the full source family, with $b,c$ fixed and $a=2s$.
Its cover is defined by
\begin{equation}
cr^3-2r^2+br-4s=0,\qquad
s=\frac{cr^3-2r^2+br}{4}.
\label{hc:full-cover}
\end{equation}
On its original inverse chart $P_r=3cr^2-4r+b=2\alpha\ne0$,
implicit differentiation gives
\begin{equation}
\partial_s^{\rm lift}=\frac4{P_r}\partial_r,\qquad
L_{b,c}=-\frac4{P_r^2}\partial_r^2
                  +\frac{4P_{rr}}{P_r^3}\partial_r.
\label{hc:full-generator}
\end{equation}
To check the second identity, apply the first operator twice:
$16P_r^{-2}\partial_r^2-16P_{rr}P_r^{-3}\partial_r$,
then multiply by $-1/4$. Thus the actual function
$Z_t((cr^3-2r^2+br)/4)$ satisfies this exact transported equation
on $P_r\ne0$. The branch equations, with no coefficient removed, are
\begin{equation}
b=4r_0-3cr_0^2,\qquad s_0=\frac{r_0^2-cr_0^3}{2}.
\label{hc:branch-family}
\end{equation}
They follow by solving $P_r=0$ and substituting in \eqref{hc:full-cover}.

For the two-sheet subfamily $c=0$, retain arbitrary $b\in\C$ and put
$\Delta=b^2-32s$. The algebraic relation is
$r^2-(b/2)r+2s=0$. In the unchanged ordered basis $(1,r)$,
\begin{equation}
R_b=\begin{pmatrix}0&-2s\\1&b/2\end{pmatrix},\qquad
D_b=\partial_s+A_b,\quad
A_b=\begin{pmatrix}0&4b/\Delta\\0&-16/\Delta\end{pmatrix}.
\label{hc:b-matrix}
\end{equation}
These formulas hold on $\Delta\ne0$. Indeed
$(b-4r)^2=\Delta$ in the algebra, hence
$\partial_s r=4/(b-4r)=4(b-4r)/\Delta$.
This gives both coefficient entries of $A_b$. Multiplication by $r$
gives $R_b$. Direct differentiation and multiplication then prove
\begin{align}
R_b^2-\tfrac b2R_b+2sI&=0,\notag\\
[D_b,R_b]&=4(bI-4R_b)/\Delta,\notag\\
D_b^2&=\partial_s^2+2A_b\partial_s+(16/\Delta)A_b.
\label{hc:b-identities}
\end{align}
For the last equality, $A_b'=(32/\Delta)A_b$ and
$A_b^2=(-16/\Delta)A_b$. At $b=0$ these identities recover
\eqref{hc:D}--\eqref{hc:D2} exactly.

The other root is $b/2-r$. Thus the sheet involution and trace in this
original basis are respectively
\[
J_b=\begin{pmatrix}1&b/2\\0&-1\end{pmatrix},\qquad
\operatorname{Tr}_b(q_0,q_1)=2q_0+(b/2)q_1.
\]
$J_b^2=I$, and $J_b$ commutes with $D_b$ because $b$ is fixed and
$A_bJ_b=J_bA_b$. Also
$\operatorname{Tr}_b D_bq=\partial_s\operatorname{Tr}_bq$, as the
two added terms are $8bq_1/\Delta-8bq_1/\Delta=0$.
The trace kernel is exactly the set of pairs $(-bq_1/4,q_1)$;
the half-trace has the section $q\mapsto(q,0)$.
This explicitly exhibits the branch information and its coupling
when the original linear coefficient is present.

The branch point is $r_0=b/4$, $s_0=b^2/32$. The exact coordinate
change $\eta=r-b/4$, with inverse $r=\eta+b/4$, gives
\begin{equation}
s=s_0-\eta^2/2,
\qquad Z_t(s)=Z_t(s_0-\eta^2/2).
\label{hc:translated-branch}
\end{equation}
The arithmetic function and its time origin have not changed.
The displayed translation applies to the cover only. For any specified
$s_0\in\C$ there is a coefficient $b$ with $b^2=32s_0$, so every
base point occurs as a branch point of an original $c=0$ fibre family.
If $Z_t(s_0)=0$ has multiplicity $m$, its factorization
$Z_t(s)=(s-s_0)^m A(s)$, $A(s_0)\ne0$, pulls back to
$(-1/2)^m\eta^{2m}A(s_0-\eta^2/2)$ and has multiplicity $2m$.
This proves the exact contact with any actual arithmetic zero and
keeps its analytic unit.

The original inverse point there has
\[
\alpha=b/2-2r=-2\eta,\quad
x=-1/(2\eta),\quad y=b/4+3\eta,\quad
w=26\eta^2+(3b/2)\eta.
\]
These follow directly from \eqref{hc:inverse}. At $\eta=0$ its two
chart branches escape, while the separate point
$(0,b,2s_0-4b^2)$ remains in the finite fibre. Thus the zero-free
origin of the $b=0$ arithmetic pullback does not remove the branch
mechanism from the entire original coefficient family. Conversely,
choosing a coefficient to place a branch at a specified base point
does not prove that that point is an arithmetic zero, or determine
whether a zero lies off the real $s$ axis. The formulas state exactly
what each operation retains.
